\documentclass[11pt]{article}
\usepackage[a4paper,margin=27mm]{geometry}
\usepackage[T1]{fontenc}
\usepackage{lmodern,microtype,amsmath,amssymb,amsthm,mathtools,enumitem}
\usepackage{xcolor}
\usepackage[colorlinks=true,linkcolor=blue,citecolor=blue,urlcolor=blue]{hyperref}
%\hypersetup{pdftitle={Coherence and Dominance under General Rankings},hypertexnames=false}
\newtheorem{theorem}{Theorem}[section]
\newtheorem{lemma}[theorem]{Lemma}
\newtheorem{proposition}[theorem]{Proposition}
\newtheorem{corollary}[theorem]{Corollary}
\theoremstyle{definition}
\newtheorem{definition}[theorem]{Definition}
\newcommand{\M}{\mathcal M}
\newcommand{\C}{\mathcal C}
\newcommand{\F}{\mathcal F}
\newcommand{\U}{\mathcal U}
\newcommand{\Lf}{\mathfrak L}
\newcommand{\supp}{\operatorname{supp}}
\setlist{itemsep=1pt,topsep=3pt}

\title{Dominance Generalised}
\author{Jason Konek}
\date{October 2026}

\begin{document}
\maketitle

\begin{abstract}
\end{abstract}


\section{Introduction}


Dominance arguments are commonly used to justify the claim that some axioms or other are requirements of rationality. For example, De~Finetti argued that an agent's forecasts or estimates of bounded random variables ought to be probabilistic expectations \cite[\S1.6, \S3.3]{deFinetti1974}. This is necessary and sufficient for any finite collection of those estimates to avoid uniform dominance when evaluated using a loss function known as the Brier score. A finite set of estimates is uniformly dominated if there is some alternative set of forecasts/estimates that is guaranteed to get a smaller penalty (according to the loss function). Dominance seems to indicate a structural defect along whatever dimension of evaluation the loss function tracks. If the Brier score tracks a dimension of evaluation that is of cardinal importance to rationality, then forecasts/estimates that are dominated are plausibly rationally impermissible. Hence non-probabilistic estimates are rationally impermissible. Probabilistic estimates, in contrast, are rationally permissible (or at least they pass every test the dominance argument considers).


A number of authors have provided dominance arguments for coherence norms like probabilism~\cite{Joyce2009,Predd2009,SSK2009,Pettigrew2022,Nielsen2023}, norms governing rational learning like conditionalization~\cite{Briggs2020,Nielsen2021},

As important as these arguments are, they also muddle together




\newpage


Many arguments for coherence share one shape. Models of some kind, such as credence functions, forecasts or sets of desirable gambles, are assessed by how well they perform in each possible world. The argument then shows two things: no coherent model is dominated, and every incoherent model is dominated. Many such arguments show more, namely that the dominating model can be taken to be coherent and itself undominated; but the conclusion that incoherent models are irrational needs only that they be dominated. De Finetti's argument from sure loss has this shape~\cite{deFinetti1974}, and so do the accuracy-dominance arguments for probabilism~\cite{Joyce2009,Predd2009,SSK2009,Pettigrew2022,Nielsen2023} and for conditionalization~\cite{Briggs2020,Nielsen2021}. The complete class theorem of statistical decision theory~\cite{Wald1950} is a relative: it identifies the undominated decision rules with Bayes rules and their limits, among randomised rules.

This paper studies such arguments in a general setting. Models carry no structure, the set of worlds is arbitrary, and a model's performance at a world is assessed by a ranking: a binary relation at each world that need not be numerical, complete or transitive. In particular, a ranking need not be generated by a numerical loss: loss rankings are an important special case, but the general theorems of Section~\ref{sec:general} use no numerical values, and the rankings constructed in Section~\ref{sec:construct} are not assumed to be transitive. The coherent models form a class fixed in advance. The argument must show that, under each ranking that evaluates it, an incoherent model is dominated by some model; Section~\ref{sec:conclusion} proves this from explicit assumptions on the rankings rather than assuming it. The main results show that every model is weakly dominated by a model from an assigned family, which yields dominance once ties are excluded. When the assigned models are undominated coherent models, they give the \emph{dominance property} of a ranking: every evaluable model is weakly dominated by an undominated coherent model. Together with closure under ties, this yields dominance of every incoherent model by an undominated coherent model. Neither dominance nor the dominance property is a standing assumption.




\newpage


\section{Setting}\label{sec:setting}

\begin{definition}\label{def:rankings}
Let $\Omega$ be a nonempty set of \emph{worlds} and $\M$ a nonempty set of \emph{models}. A \emph{ranking} $L$ consists of a nonempty set $\M_L\subseteq\M$, whose members are the models \emph{evaluable} by $L$, and, for each world $\omega$, a binary relation $\preceq_\omega$ on $\M_L$. We read $m\preceq_\omega m'$ as: $m$ performs at least as well as $m'$ at $\omega$. Models $m$ and $m'$ are \emph{tied at $\omega$} if $m\preceq_\omega m'$ and $m'\preceq_\omega m$. No property of the relations is assumed. $L$ is \emph{transitive} if every $\preceq_\omega$ is transitive. For a fixed ranking $L$, ``model'' means a member of $\M_L$.
\end{definition}

\begin{definition}\label{def:dominance}
Fix a ranking.
\begin{itemize}[leftmargin=*]
\item $c$ \emph{weakly dominates} $m$ if $c\preceq_\omega m$ at every world.
\item $m$ and $c$ are \emph{tied at every world} if each weakly dominates the other.
\item $c$ \emph{dominates} $m$ if $c$ weakly dominates $m$ and they are not tied at every world.
\end{itemize}
A model is \emph{undominated} if no model dominates it. $\U_L$ is the set of undominated models.
\end{definition}

\begin{definition}\label{def:priors}
For a finite nonempty $F\subseteq\Omega$, let $\Delta_F$ be the set of probability mass functions (pmfs) on $F$, $p=(p_\omega)_{\omega\in F}$, with the Euclidean topology. The support of a pmf $p$ is $\supp p\coloneqq\{\omega:p_\omega>0\}$. Let $\delta_\omega$ be the point mass at $\omega$. When $\Omega$ is finite, $\Delta\coloneqq\Delta_\Omega$.
\end{definition}

\begin{definition}\label{def:coherence-property}
A class $\C\subseteq\M$ of \emph{coherent} models is fixed in advance, independently of the rankings. For a ranking $L$, the \emph{undominated coherent models} of $L$ are
\[
\C_L\coloneqq\{c\in\C\cap\M_L:\ \text{no } m\in\M_L \text{ dominates } c\}=\C\cap\U_L.
\]
%$L$ has the \emph{dominance property} if every model in $\M_L$ is weakly dominated by a member of $\C_L$.
\end{definition}

\subsection{Axioms}

\begin{definition}\label{def:tie-closure}
A set $\mathcal A$ of models is \emph{closed under ties} if every model tied at every world with a member of $\mathcal A$ belongs to $\mathcal A$. A ranking $L$ is \emph{closed under ties} if $\C_L$ is.
\end{definition}

\begin{definition}\label{def:family-assumptions}
Let $\Lf$ be a family of rankings.
\begin{itemize}[leftmargin=*]
\item \emph{Evaluability}: every model in $\M$ is evaluable by some ranking in $\Lf$.
\item \emph{Coherent admissibility}: every coherent model belongs to $\C_L$ for some $L\in\Lf$.
\end{itemize}
\end{definition}

\begin{lemma}\label{lem:ties}
Let a model $a$ weakly dominate a model $m$. Then $a$ dominates $m$ if and only if $m$ does not weakly dominate $a$.
\begin{enumerate}[label=(\alph*)]
\item If $m$ is undominated, then $m$ and $a$ are tied at every world.
\item If $L$ is closed under ties, $a\in\C_L$ and $m\notin\C_L$, then $a$ dominates $m$.
\item More generally, if $\mathcal A$ is closed under ties, $a\in\mathcal A$ and $m\notin\mathcal A$, then $a$ dominates $m$.
\end{enumerate}
\end{lemma}
\begin{proof}
The first claim is the definition of dominance. (a) Otherwise $a$ would dominate $m$. (c) Since $m\notin\mathcal A$, $m$ is not tied at every world with $a$. (b) This is (c) with $\mathcal A\coloneqq\C_L$.
\end{proof}
Undominatedness of $a$ cannot replace closure under ties in (c): with one world and two models $a$ and $m$, each at least as good as the other, neither is dominated, and $a$ weakly dominates $m$ without dominating it.

\section{The dominance theorems}\label{sec:general}

\begin{lemma}\label{lem:kkm}
Let $F$ be finite and nonempty, and let $(A_\omega)_{\omega\in F}$ be closed subsets of $\Delta_F$ such that every $p\in\Delta_F$ lies in $A_\omega$ for some $\omega\in\supp p$. Then $\bigcap_{\omega\in F}A_\omega\ne\varnothing$.
\end{lemma}
\begin{proof}
Each $A_\omega$ is nonempty, since by hypothesis we always have 
\[
\delta_\omega\in\bigcup_{\omega^\prime\in\supp \delta_\omega}A_{\omega^\prime}=A_\omega.
\]
Let $d_\omega(p)\coloneqq\min_{q\in A_\omega}\lVert p-q \rVert$ and $s(p)\coloneqq\sum_{\omega\in F} d_\omega(p)$. The map $G(p)_\omega\coloneqq(p_\omega+d_\omega(p))/(1+s(p))$ is continuous. Moreover, $G$ maps each $p\in\Delta_F$ to some point in $\Delta_F$ (since each coordinate $G(p)_\omega$ of $G(p)$ is nonnegative and these coordinates sum to 1). Brouwer's fixed-point theorem then guarantees that there is some $p\in\Delta_F$ with $G(p)=p$. This holds just in case $p_\omega\cdot s(p)=d_\omega(p)$ for all $\omega\in F$. By assumption, there is some $\omega\in\supp p$ with $p\in A_\omega$. Choose some such $\omega$. Then $d_\omega(p)=0$, so $p_\omega\cdot s(p)=0$. Since $p_\omega>0$, $s(p)=0$ and hence $p\in\bigcap_{\omega\in F}A_\omega$.
\end{proof}

This is the lemma of Knaster, Kuratowski and Mazurkiewicz~\cite{KKM1929}.




\subsection{Finitely many worlds}
\begin{theorem}\label{thm:finite}
Let $\Omega$ be finite and let $L$ be any ranking. Suppose that for each pmf $p\in\Delta$, there is an assigned model $c_p$, not necessarily coherent or undominated. And for each model $m$ and world $\omega$, there is a closed set $A_\omega(m)\subseteq\Delta$, such that:
\begin{enumerate}[label=(\roman*)]
\item $c_p\preceq_\omega m$ for every $p\in A_\omega(m)$;
\item for every model $m$ and every pmf $p\in\Delta$, there is some $\omega\in\supp p$ such that $p\in A_\omega(m)$.
\end{enumerate}
Then for each model $m$, some pmf lies in every $A_\omega(m)$, and the model assigned to any such pmf weakly dominates $m$.
\end{theorem}
\begin{proof}
Choose $m$. By (ii), the closed sets $A_\omega(m)$ satisfy the hypothesis of Lemma~\ref{lem:kkm}, so some $p^*$ lies in all of them. By (i), $c_{p^*}\preceq_\omega m$ at every world.
\end{proof}

\begin{corollary}\label{cor:closed-choice}
Let $\Omega$ be finite. Suppose that for each $p\in\Delta$ there is an assigned model $c_p$, not necessarily coherent or undominated, such that every model $m$ has some $\omega\in\supp p$ with $c_p\preceq_\omega m$. Suppose further that for any model $m$ and world $\omega$, the set $\{p:c_p\preceq_\omega m\}$ is closed. Then every model is weakly dominated by some $c_p$.
\end{corollary}
\begin{proof}
Take $A_\omega(m)\coloneqq\{p:c_p\preceq_\omega m\}$. Then condition (i) of Theorem~\ref{thm:finite} holds by definition of $A_\omega(m)$, and condition (ii) holds by the first assumption.
\end{proof}

\subsection{Arbitrary sets of worlds}
%A space is \emph{compact} if every open cover has a finite subcover; no separation property is assumed.

\begin{theorem}\label{thm:local}
Let $\Omega$ be any nonempty set and $L$ any ranking. Suppose:
\begin{enumerate}[label=\arabic*.,leftmargin=*]
\item There is a compact space of parameters $\Theta$. For each $\theta\in\Theta$, there is a selected model $c_\theta$, not necessarily coherent or undominated. For each world $\omega$ and model $m$ there is a closed set $K_\omega(m)\subseteq\Theta$ such that $c_\theta\preceq_\omega m$ for every $\theta\in K_\omega(m)$.
\item For every finite set of worlds $G\subseteq\Omega$, there are a nonempty finite set $F\supseteq G$, a map $\theta_F:\Delta_F\to\Theta$, and, for each $\omega\in F$ and model $m$, a closed set $A^F_\omega(m)\subseteq\Delta_F$, such that
\begin{enumerate}[label=(\alph*)]
\item for every $\omega\in F$ and model $m$, $\theta_F(p)\in K_\omega(m)$ for all $p\in A^F_\omega(m)$;
\item for every model $m$ and every pmf $p\in\Delta_F$, there is some $\omega\in\supp p$ such that $p\in A^F_\omega(m)$.
\end{enumerate}
\end{enumerate}
Then every model is weakly dominated by some $c_\theta$.
\end{theorem}
\begin{proof}
Choose a model $m$. Let $F$ be any set provided by assumption~2. The closed sets $A^F_\omega(m)$, $\omega\in F$, satisfy the hypothesis of Lemma~\ref{lem:kkm}, by assumption~2(b). So some $p_F$ lies in their intersection. By assumption~2(a), $\theta_F(p_F)\in\bigcap_{\omega\in F}K_\omega(m)$. For any finite number of worlds $\omega_1,\dots,\omega_k$, apply assumption~2 to $G\coloneqq\{\omega_1,\dots,\omega_k\}$ to obtain $F\supseteq G$. Then $\bigcap_iK_{\omega_i}(m)\supseteq\bigcap_{\omega\in F}K_\omega(m)\ni\theta_F(p_F)$. So every finite subfamily of the closed sets $K_\omega(m)$, $\omega\in\Omega$, has a nonempty intersection. It follows that $\bigcap_{\omega\in\Omega}K_\omega(m)\ne\varnothing$. Otherwise the open sets $\Theta\setminus K_\omega(m)$, $\omega\in\Omega$, would cover $\Theta$. By compactness, finitely many of them, $\Theta\setminus K_{\omega_1}(m),\dots,\Theta\setminus K_{\omega_k}(m)$, would cover $\Theta$. But then $K_{\omega_1}(m)\cap\dots\cap K_{\omega_k}(m)=\varnothing$, a contradiction. 

Choose $\theta^*\in\bigcap_{\omega\in\Omega}K_\omega(m)$. By assumption~1, $c_{\theta^*}\preceq_\omega m$ for every world $\omega$. That is, $c_{\theta^*}$ weakly dominates $m$.
\end{proof}

None of these proofs uses coherence or undominatedness of the assigned models. Weak domination is not yet dominance: an assigned model may be tied at every world with the model it weakly dominates. The following corollary states what excludes such ties.

\begin{corollary}\label{cor:coherent-assigned}
Assume the hypotheses of Theorem~\ref{thm:finite}, Corollary~\ref{cor:closed-choice} or Theorem~\ref{thm:local}, and let $\mathcal A$ be the set of assigned models.
\begin{enumerate}[label=(\alph*)]
\item If $\mathcal A$ is closed under ties, every model not in $\mathcal A$ is dominated by a member of $\mathcal A$.
\item If $\mathcal A\subseteq\C_L$, every undominated model is tied at every world with a member of $\C_L$.
\item If $\mathcal A\subseteq\C_L$ and $L$ is closed under ties, then $\U_L=\C_L$, and every other model, in particular every incoherent model, is dominated by a member of $\C_L$.
\end{enumerate}
\end{corollary}
\begin{proof}
Each model $m$ is weakly dominated by some $a\in\mathcal A$. (a) is Lemma~\ref{lem:ties}(c). (b) If $m$ is undominated, Lemma~\ref{lem:ties}(a) ties it to $a\in\C_L$. (c) By (b) and closure under ties, $\U_L\subseteq\C_L\subseteq\U_L$. If $m\notin\C_L$, then $a$ dominates $m$ by Lemma~\ref{lem:ties}(b). Incoherent models are not in $\C_L\subseteq\C$.
\end{proof}

\section{The dominance argument}\label{sec:conclusion}

The argument has three ingredients: a dominance principle, which says when a model is irrational and when it is rationally permissible (Definition~\ref{def:permissibility}); a list of every assumption the argument makes about the rankings (Definition~\ref{def:argument-assumptions}); and a theorem that derives from these assumptions that the incoherent models are irrational and the coherent models rationally permissible (Theorem~\ref{thm:argument}).

\begin{definition}[Dominance principle]\label{def:permissibility}
Let $\Lf$ be a family of rankings. A model is \emph{irrational} if, under every ranking in $\Lf$ that evaluates it, some model dominates it. A model is \emph{rationally permissible} if it is undominated under some ranking in $\Lf$ that evaluates it.
\end{definition}

\begin{lemma}\label{lem:complement}
A model is irrational if and only if it is not rationally permissible. Under evaluability, an irrational model is dominated under at least one ranking in $\Lf$.
\end{lemma}
\begin{proof}
A model fails to be irrational exactly when some ranking in $\Lf$ that evaluates it leaves it undominated, which is rational permissibility; by evaluability, some ranking in $\Lf$ evaluates it.
\end{proof}
The equivalence uses no transitivity, compactness or undominated dominators.

\begin{definition}[Assumptions on the rankings]\label{def:argument-assumptions}
The dominance argument assumes the following of the family $\Lf$.
\begin{itemize}[leftmargin=2.4em,labelsep=0.5em]
\item[\textbf{(E)}] \emph{Evaluability}: every model in $\M$ is evaluable by some ranking in $\Lf$. No model escapes assessment.
\item[\textbf{(A)}] \emph{Coherent admissibility}: every coherent model belongs to $\C_L$ for some $L\in\Lf$. Each coherent model is undominated under at least one ranking that evaluates it.
\item[\textbf{(S)}] \emph{Selection}: for each $L\in\Lf$, the hypotheses of Theorem~\ref{thm:finite}, of Corollary~\ref{cor:closed-choice} or of Theorem~\ref{thm:local} hold, with every assigned model $c_p$, or selected model $c_\theta$, in $\C_L$. Each ranking has enough undominated coherent models to meet the conditions of one of the theorems of Section~\ref{sec:general}.
\item[\textbf{(T)}] \emph{Closure under ties}: every $L\in\Lf$ is closed under ties, that is, $\C_L$ is closed under ties. A model tied at every world with an undominated coherent model is itself an undominated coherent model.
\end{itemize}
\end{definition}

Conditions (E) and (A) are those of Definition~\ref{def:family-assumptions}, and (T) is that of Definition~\ref{def:tie-closure}. Beyond the setting of Definition~\ref{def:rankings}, these are all the assumptions about the rankings that the argument uses.

Three things are not assumed. First, no property of the relations $\preceq_\omega$ is assumed: the rankings need not be transitive or complete. Second, it is not assumed that any incoherent model is dominated. Condition (A) concerns coherent models only. Condition (S) concerns the assigned models: they must be undominated coherent models and must meet the conditions of one of the theorems of Section~\ref{sec:general}, which compare them with each model at particular worlds; it says nothing about which other models are dominated. Condition (T) concerns ties: it implies that no incoherent model is tied at every world with an undominated coherent model, but not that any incoherent model is dominated. That every incoherent model is dominated, and so irrational, is proved in Theorem~\ref{thm:argument}. Third, no ranking is required to leave every coherent model undominated.

Closure under ties cannot be dropped. Take one world, a coherent model $a$ and an incoherent model $m$, with every comparison holding, and let $\Lf$ consist of this one ranking. Then $\C_L=\{a\}$, and (E), (A) and (S) hold, with $a$ assigned to the only pmf. But $m$ is undominated, so it is rationally permissible.

\begin{theorem}[The dominance argument]\label{thm:argument}
Let $\Lf$ satisfy (E), (A), (S) and (T). Then the rationally permissible models are exactly the coherent ones, and the irrational models are exactly the incoherent ones. Moreover, for each $L\in\Lf$, every incoherent model in $\M_L$ is dominated by a member of $\C_L$, and $\U_L=\C_L$. Every incoherent model is dominated under at least one ranking in $\Lf$.
\end{theorem}
\begin{proof}
\leavevmode
\begin{enumerate}[label=\arabic*.,leftmargin=*]
\item Fix $L\in\Lf$ and an incoherent model $m\in\M_L$. By (S), Theorem~\ref{thm:finite}, Corollary~\ref{cor:closed-choice} or Theorem~\ref{thm:local} gives an assigned model $a$ that weakly dominates $m$, and $a\in\C_L$.
\item Since $\C_L\subseteq\C$ and $m$ is incoherent, $m\notin\C_L$. By (T), $L$ is closed under ties. So $a$ dominates $m$, by Lemma~\ref{lem:ties}(b).
\item Since $L$ was arbitrary, every incoherent model is dominated under every ranking in $\Lf$ that evaluates it; that is, it is irrational. By Lemma~\ref{lem:complement}, it is not rationally permissible, and, by (E), it is dominated under at least one ranking in $\Lf$.
\item By (A), each coherent model $c$ belongs to $\C_L$ for some $L\in\Lf$. Then $L$ evaluates $c$ and leaves it undominated, so $c$ is rationally permissible. By Lemma~\ref{lem:complement}, it is not irrational.
\item Every model is coherent or incoherent. So, by steps 3 and 4, the rationally permissible models are exactly the coherent ones, and the irrational models are exactly the incoherent ones.
\item Finally, for each $L\in\Lf$, (S) and (T) are the hypotheses of Corollary~\ref{cor:coherent-assigned}(c), which gives $\U_L=\C_L$.\qedhere
\end{enumerate}
\end{proof}

Condition (E) is used only for the last claim, so that a verdict of irrationality is not vacuous: without (E), a model that no ranking evaluates would count as irrational trivially. Steps 1 and 6 use (S) only through its consequence that every model in $\M_L$ is weakly dominated by a member of $\C_L$.

In summary, the argument runs as follows.
\begin{enumerate}[label=\textbf{(P\arabic*)},leftmargin=*]
\item The family $\Lf$ satisfies (E), (A), (S) and (T).
\item \emph{Dominance principle} (Definition~\ref{def:permissibility}). A model is irrational exactly when it is dominated under every ranking in $\Lf$ that evaluates it, and rationally permissible exactly when it is undominated under some ranking in $\Lf$ that evaluates it; so no model is both (Lemma~\ref{lem:complement}).
\item Theorem~\ref{thm:argument}. Given (P1), every incoherent model is dominated under every ranking in $\Lf$ that evaluates it, and every coherent model is undominated under some ranking in $\Lf$ that evaluates it.
\item[\textbf{(C)}] The rationally permissible models are exactly the coherent ones, and the irrational models are exactly the incoherent ones.
\end{enumerate}
Premise (P1) must be verified for the rankings at hand, (P2) is the normative principle, and (P3) is proved.

For sets of desirable gambles, the axioms (A1)--(A4) of~\cite{Konek2026} can replace (S) and (T): for every family of rankings satisfying them, given a measure satisfying its conditions (M2) and (M3$'$), every incoherent model is dominated under every ranking that evaluates it~\cite[Theorem~4.1]{Konek2026}, so steps 3--5 of the proof apply unchanged, although that theorem gives ordinary dominance, not undominated coherent dominators.

\subsection{Further conclusions under extra hypotheses}
The proof of Theorem~\ref{thm:argument} uses (S) only through its consequence that every model is weakly dominated by a member of $\C_L$. The following results take that consequence as their hypothesis, in place of (S), and add closure under ties or transitivity. This hypothesis alone does not make incoherent models dominated, since an incoherent model may be tied at every world with its weak dominator. Under closure under ties it does, with dominators in $\C_L$. Transitivity alone does not: an incoherent model tied at every world with a member of $\C_L$ is then undominated (Proposition~\ref{prop:profiles}). Under transitivity, the models that are dominated have dominators in $\C_L$.

\begin{proposition}\label{prop:conclusion}
Let $L$ be a ranking under which every model is weakly dominated by a member of $\C_L$. Then every undominated model is tied at every world with a member of $\C_L$. If $L$ is also closed under ties, then $\U_L=\C_L$, and every other model is dominated by a member of $\C_L$.
\end{proposition}
\begin{proof}
Apply Lemma~\ref{lem:ties} to each model and a member of $\C_L$ weakly dominating it. By part (a) of Lemma~\ref{lem:ties}, every undominated model $m$ is tied with a member of $\C_L$. If $L$ is also closed under ties, then $m\in\C_L$. Hence $\U_L=\C_L$. Part (b) of Lemma~\ref{lem:ties} gives the last claim.
\end{proof}

\begin{proposition}\label{prop:profiles}
Let $L$ be a transitive ranking under which every model is weakly dominated by a member of $\C_L$. Then $\U_L$ consists of the models tied at every world with some member of $\C_L$, and every other model is dominated by a member of $\C_L$.
\end{proposition}
\begin{proof}
By Proposition~\ref{prop:conclusion}, every undominated model is tied with some $c\in\C_L$. Conversely, let $m$ be tied with $c\in\C_L$, and suppose $m'$ dominates $m$. Then $m'\preceq_\omega m\preceq_\omega c$ at every world, so $m'$ weakly dominates $c$ by transitivity. If $c$ weakly dominated $m'$, then $m\preceq_\omega c\preceq_\omega m'$ would make $m$ weakly dominate $m'$, which is false. So $m'$ dominates $c$, which is impossible. Hence $m$ is undominated. Finally, if $m$ is dominated, let $c\in\C_L$ weakly dominate it. Since $m$ is not tied with $c$, $c$ dominates $m$.
\end{proof}

\begin{proposition}\label{prop:permissible}
Suppose that under each ranking $L\in\Lf$, every model in $\M_L$ is weakly dominated by a member of $\C_L$. Then for each $L\in\Lf$:
\begin{enumerate}[label=(\alph*)]
\item every undominated model is tied at every world with a member of $\C_L$;
\item if $L$ is closed under ties, then $\U_L=\C_L$, and every other model in $\M_L$ is dominated by a member of $\C_L$;
\item if $L$ is transitive, then $\U_L$ consists of the models tied at every world with a member of $\C_L$, and every other model in $\M_L$ is dominated by a member of $\C_L$.
\end{enumerate}
If every ranking in $\Lf$ is closed under ties or transitive, then a model is irrational iff, under every ranking in $\Lf$ that evaluates it, some member of $\C_L$ dominates it. If every ranking in $\Lf$ is closed under ties and coherent admissibility holds, the rationally permissible models are exactly the coherent ones.
\end{proposition}
\begin{proof}
Parts (a) and (b) are Proposition~\ref{prop:conclusion}, and (c) is Proposition~\ref{prop:profiles}. If $m$ is irrational and $L$ evaluates $m$, then $m\notin\U_L$, so by (b) or (c) some member of $\C_L$ dominates $m$. The converse is immediate from the definitions. If every ranking is closed under ties, (b) gives $\U_L=\C_L\subseteq\C$ for every $L$, so every rationally permissible model is coherent; coherent admissibility makes every coherent model rationally permissible.
\end{proof}
The last claim is the main conclusion of Theorem~\ref{thm:argument}, with (S) replaced by the hypothesis of this subsection; it does not need (E). What Proposition~\ref{prop:permissible} adds is the case of transitive rankings that are not closed under ties, in which irrational models still have dominators in $\C_L$.

% Insert before the bibliography of the Short paper.
% Add the entries in Applications-bibliography.tex to its bibliography.
\section{Applications to accuracy}\label{sec:short-applications}
For a loss $L:\M_L\times\Omega\to[0,\infty]$, put $m\preceq_\omega m'$ exactly when $L(m,\omega)\le L(m',\omega)$. A \emph{Bayes model for $p$} minimises expected loss, with $0\cdot\infty=0$. The conclusions below concern dominance as defined here: improvement at some world, without worsening at any. Improvement at every world is a stronger conclusion. {\color{red} For a finitely supported prior $p$, write $E_pL(m)\coloneqq\sum_{\omega\in\supp p}p_\omega L(m,\omega)$. Each proof below verifies selection (S) and closure under ties (T) for the ranking at hand, and then applies Corollary~\ref{cor:coherent-assigned}(c).}

\subsection{Proper scores}
\begin{proposition}\label{prop:short-proper}
{\color{red}Let $\Omega$ be finite and $\C\subseteq\M_L$. Suppose that $p\mapsto c_p$ maps $\Delta$ onto $\C$, that each $c_p$ is the unique Bayes model for $p$, and that $p\mapsto L(c_p,\omega)$ is lower semicontinuous for every world $\omega$. Then the undominated models are exactly the coherent models, and every other model is dominated by a coherent model. In particular, if $\C=\Delta$, each $p\in\Delta$ is the unique Bayes model for $p$, and $p\mapsto L(p,\omega)$ is lower semicontinuous for every world, then the undominated models are exactly the probabilities, and every other model is dominated by a probability.}
\end{proposition}
\begin{proof}
{\color{red}If $E_pL(c_p)=\infty$ for some $p$, every model is a Bayes model for $p$, so $\M_L=\{c_p\}$ and the claims are trivial. Assume otherwise.

\emph{Closure under ties.} If $m$ weakly dominates $c_p$, then $E_pL(m)\le E_pL(c_p)$, so $m$ is a Bayes model for $p$, and $m=c_p$. So every coherent model is undominated, $\C_L=\C$, and the only model tied at every world with $c_p$ is $c_p$ itself: (T) holds.

\emph{Selection.} If some $m$ had $L(m,\omega)<L(c_p,\omega)$ at every $\omega\in\supp p$, then $E_pL(m)<E_pL(c_p)$, as $E_pL(c_p)$ is finite; this contradicts Bayes optimality. So every model $m$ has some $\omega\in\supp p$ with $c_p\preceq_\omega m$. Each set $\{p:L(c_p,\omega)\le L(m,\omega)\}$ is $\Delta$ if $L(m,\omega)=\infty$, and closed by lower semicontinuity otherwise. So the hypotheses of Corollary~\ref{cor:closed-choice} hold with every $c_p\in\C_L$: (S) holds. Corollary~\ref{cor:coherent-assigned}(c) completes the proof.}
\end{proof}
This includes strictly proper scores that are not additive, as in Pettigrew~\cite{Pettigrew2022}, and permits infinite losses. {\color{red}It also applies to finitely many event forecasts on arbitrary $\Omega$: replace worlds by their finitely many truth-value patterns $w$, with truth-value vectors $v_w$, and put $c_p\coloneqq\sum_wp_wv_w$; the coherent forecasts are exactly these vectors. Continuous, strictly proper component scores make $c_p$ the unique Bayes forecast for $p$, recovering} the ordinary-dominance consequence of Predd et al.~\cite{Predd2009}. Those results also establish stronger forms of dominance. For merely proper scores, uniqueness can fail: incoherent forecasts can share coherent loss profiles and remain undominated~\cite{SSK2009}{\color{red}; then (T) fails}. This is precisely the distinction made by closure under ties.

\subsection{Probabilism on arbitrary domains}
Let $\M_L=[0,1]^{2^\Omega}$ and let $\C$ be the finitely additive probabilities on all subsets of $\Omega$, with the topology of pointwise convergence. For $f\ge0$, define $E_qf=\sup\{E_qh:0\le h\le f,\ h\text{ bounded}\}$. Call $L$ \emph{quasi-strictly proper} if
\[
E_qL(q,\cdot)\le E_qL(m,\cdot)\quad(q\in\C,\ m\in\M_L),
\]
strictly whenever $m\notin\C$.
\begin{proposition}\label{prop:short-infinite}
If $L$ is quasi-strictly proper and each $q\mapsto L(q,\omega)$ is lower semicontinuous, the undominated models are exactly the finitely additive probabilities, and every other model is dominated by one.
\end{proposition}
\begin{proof}
First suppose $m$ dominates $q\in\C$, strictly improving at $w$. Quasi-strict propriety gives $E_qL(m)=E_qL(q)<\infty$; finiteness follows by comparison with any incoherent credence. Put $q_n=(1-1/n)q+\delta_w/n$. Propriety at $q_n$ and $q$ gives
\[
\tfrac1nL(q_n,w)+(1-\tfrac1n)E_qL(q)
\le E_{q_n}L(q_n)\le E_{q_n}L(m)
=\tfrac1nL(m,w)+(1-\tfrac1n)E_qL(q).
\]
Hence $L(q_n,w)\le L(m,w)$, contradicting lower semicontinuity at $q$. So probabilities are undominated{\color{red}, and $\C_L=\C$. If an incoherent $m$ were tied at every world with $q\in\C$, then $E_qL(m,\cdot)=E_qL(q,\cdot)$, contradicting quasi-strict propriety. So (T) holds.}

{\color{red}For (S), apply Theorem~\ref{thm:local} with the compact space $\Theta\coloneqq\C$, which is closed in $[0,1]^{2^\Omega}$, $c_q\coloneqq q\in\C_L$, and $K_\omega(m)\coloneqq\{q:L(q,\omega)\le L(m,\omega)\}$, which is closed by lower semicontinuity. Given a finite set $G$ of worlds, let $F\coloneqq G$, or a single world if $G=\varnothing$, let $\theta_F(p)\coloneqq p$, viewed as a finitely supported probability, which is continuous, and let $A^F_\omega(m)\coloneqq\theta_F^{-1}(K_\omega(m))$, which is closed. Then condition~2(a) holds by definition. For condition~2(b), let $p\in\Delta_F$. Quasi-strict propriety gives $E_pL(p,\cdot)<E_pL(m',\cdot)$ for an incoherent $m'$, such as the zero credence function, so $E_pL(p,\cdot)$ is finite, and $E_pL(p,\cdot)\le E_pL(m,\cdot)$. If $L(m,\omega)<L(p,\omega)$ at every $\omega\in\supp p$, then $E_pL(m,\cdot)<E_pL(p,\cdot)$, a contradiction. So (S) holds, and Corollary~\ref{cor:coherent-assigned}(c) completes the proof.}
\end{proof}
Nielsen~\cite{Nielsen2023} assumes continuity of $q\mapsto E_pL(q)$ for every probability $p$, together with local boundedness of coherent loss profiles, and obtains improvement at every world. The proposition requires only worldwise lower semicontinuity and drops boundedness, but concludes ordinary dominance. It does not claim Nielsen's stronger conclusion under weaker hypotheses.


\subsection{Conditional updating}
Let $\Omega$ be finite and $E_1,\ldots,E_k$ a partition into nonempty cells. A strategy $(c,a_1,\ldots,a_k)$ specifies current credences and credences to adopt on learning each cell. Call it coherent when all these credences are probabilities,
\[
a_j(E_j)=1,\qquad c(A\cap E_j)=c(E_j)a_j(A)\quad\text{for every }A,j.
\]
Thus even a plan for a probability-zero cell must be probabilistic and concentrate on that cell. Let $J$ be a finite loss for credence functions on $2^\Omega$, strictly proper in the sense that each probability uniquely minimises its own expected loss, and lower semicontinuous on probabilities at every world, and assess strategies by
\[
L((c,a),\omega)=\alpha J(c,\omega)+\beta J(a_j,\omega)
\quad(\omega\in E_j),\qquad \alpha,\beta>0.
\]
\begin{proposition}\label{prop:short-updating}
The undominated strategies are exactly the coherent ones, and every incoherent strategy is dominated by a coherent one.
\end{proposition}
\begin{proof}
A weak dominator of a coherent strategy first minimises its $c$-expected loss, forcing its current credences to equal $c$ by strict propriety. Cancelling that common loss and taking expectations under each $a_j$ forces every conditional report to equal $a_j$. So coherent strategies are undominated{\color{red}, $\C_L=\C$, and (T) holds}.

For selection, put $\eta=1/(2k)$ and fix a probability $u_j$ on each cell. Given $p\in\Delta$, set
\[
r_j=\frac{p|_{E_j}+(\eta-p(E_j))_+u_j}{\max\{p(E_j),\eta\}},\qquad
\lambda_j=\frac{(p(E_j)-\eta)_+}{\sum_i(p(E_i)-\eta)_+},\qquad
q=\sum_j\lambda_jr_j,
\]
where $x_+=\max\{x,0\}$. These define a continuous coherent strategy $(q,r_1,\ldots,r_k)$, Bayes for $q$, with $\supp q\subseteq\supp p$. Its Bayes inequality therefore supplies a defending world in $\supp p$ against each rival. Its losses are lower semicontinuous{\color{red} in $p$, since $J$ is lower semicontinuous on probabilities and $q$ and the $r_j$ are continuous in $p$. So the hypotheses of Corollary~\ref{cor:closed-choice} hold with every assigned strategy in $\C_L$: (S) holds, and Corollary~\ref{cor:coherent-assigned}(c) completes the proof}.
\end{proof}
The selected strategy need not be Bayes for the indexing prior $p$: a Bayes prior supported within $\supp p$ suffices. This avoids discontinuity of conditional probabilities at null cells. The result extends the finite-valued setting of Nielsen~\cite{Nielsen2021} beyond additive, extensional, continuous scores. It also isolates the null-cell issue in Briggs and Pettigrew~\cite{Briggs2020}: merely conditionalizing on positive-probability cells does not ensure undominatedness. Invalid plans on null cells can be dominated without improvement at every world. Finiteness of $J$ matters: logarithmic losses can make such plans tied with coherent strategies, violating closure under ties.

\subsection{Why coherent admissibility is insufficient}
Joyce's~\cite{Joyce2009} truth-directedness requires loss to decrease whenever every coordinate of a credence moves towards its truth value, with at least one strict improvement. Together with continuity, symmetry and coherent admissibility, this still does not ensure dominance of incoherent credences. For two worlds, take $\M_L=[0,1]^2$, $\C=\Delta$ and
\[
L_1(x)=(1-x_1+x_2)^2,\qquad L_2(x)=(1+x_1-x_2)^2.
\]
These losses have all four properties. Writing $d=x_1-x_2$, every $x$ is tied with the probability $((1+d)/2,(1-d)/2)$; distinct profiles are incomparable. Thus no model is dominated. {\color{red}Here (S) holds: $E_pL(x)=p_1(1-d)^2+p_2(1+d)^2$ is minimised at $d=p_1-p_2$, so each $p$ is a Bayes model for itself, hence competitive on $\supp p$; the losses are continuous; and $\C_L=\Delta$, since no model is dominated. So Corollary~\ref{cor:closed-choice} applies with $c_p=p$. But (T) fails: every incoherent $x$ is tied at every world with a probability.} Coherent admissibility alone cannot exclude incoherent models that perform exactly as well as coherent ones.


{\color{red}
\subsection{Chance dominance}\label{sec:chance}
Pettigrew's~\cite[Chapter~10]{Pettigrew2016} chance-dominance argument compares expected rather than actual accuracy. Let $W$ be a finite nonempty set of worlds, $\mathcal F$ a set of propositions (subsets of $W$), and $\M_L\coloneqq[0,1]^{\mathcal F}$ the credence functions. Each probability $\theta$ on $W$ induces the credence function $X\mapsto\sum_{w\in X}\theta(w)$, also written $\theta$. Let $\mathcal K$ be a nonempty set of probabilities on $W$: the possible current chance functions, conditional on the agent's total evidence. Let $\C$ consist of the credence functions induced by the closed convex hull $\operatorname{cl}\operatorname{conv}\mathcal K$; these are the credence functions satisfying Probabilism and Pettigrew's ETP$^+$. Let $I:\M_L\times W\to[0,\infty)$ be a finite inaccuracy measure such that, for every probability $\theta$ on $W$, $\theta$ is the unique minimiser of $c\mapsto E_\theta I(c)\coloneqq\sum_w\theta(w)I(c,w)$ over $\M_L$, and $c\mapsto I(c,w)$ is lower semicontinuous on the credence functions induced by probabilities, for every $w$. The worlds of the ranking are the members of $\mathcal K$, and $c\preceq_{\mathrm{ch}}d$ exactly when $E_{\mathrm{ch}}I(c)\le E_{\mathrm{ch}}I(d)$.
\begin{proposition}\label{prop:short-chance}
The undominated credence functions are exactly the members of $\C$, and every other credence function is dominated by a member of $\C$.
\end{proposition}
\begin{proof}
\emph{Closure under ties.} Let $\theta\in\operatorname{cl}\operatorname{conv}\mathcal K$, and let $m$ weakly dominate $\theta$: $E_{\mathrm{ch}}I(m)\le E_{\mathrm{ch}}I(\theta)$ for every $\mathrm{ch}\in\mathcal K$. Since $W$ is finite and $I$ is finite, $\rho\mapsto E_\rho I(m)-E_\rho I(\theta)$ is affine and continuous on the probabilities on $W$. It is nonpositive on $\mathcal K$, hence on $\operatorname{cl}\operatorname{conv}\mathcal K$, and in particular at $\rho=\theta$. So $E_\theta I(m)\le E_\theta I(\theta)$, and $m=\theta$ by strict propriety. Hence every member of $\C$ is undominated, $\C_L=\C$, and the only model tied at every world with a member $\theta$ of $\C$ is $\theta$: (T) holds.

\emph{Selection.} Apply Theorem~\ref{thm:local} with the compact space $\Theta\coloneqq\operatorname{cl}\operatorname{conv}\mathcal K$, $c_\theta\coloneqq\theta\in\C_L$, and $K_{\mathrm{ch}}(m)\coloneqq\{\theta\in\Theta:E_{\mathrm{ch}}I(\theta)\le E_{\mathrm{ch}}I(m)\}$, which is closed since $\theta\mapsto E_{\mathrm{ch}}I(\theta)=\sum_w\mathrm{ch}(w)I(\theta,w)$ is lower semicontinuous. Given a finite $G\subseteq\mathcal K$, let $F\coloneqq G$, or a single member of $\mathcal K$ if $G=\varnothing$; let $\theta_F(p)\coloneqq\sum_{\mathrm{ch}\in F}p_{\mathrm{ch}}\,\mathrm{ch}$, which is continuous; and let $A^F_{\mathrm{ch}}(m)\coloneqq\theta_F^{-1}(K_{\mathrm{ch}}(m))$, which is closed. Condition~2(a) holds by definition. For condition~2(b), note that $\sum_{\mathrm{ch}\in F}p_{\mathrm{ch}}E_{\mathrm{ch}}I(c)=E_{\theta_F(p)}I(c)$ for every $c$, which is minimised at $c=\theta_F(p)$. If $E_{\mathrm{ch}}I(m)<E_{\mathrm{ch}}I(\theta_F(p))$ for every $\mathrm{ch}\in\supp p$, then $E_{\theta_F(p)}I(m)<E_{\theta_F(p)}I(\theta_F(p))$, a contradiction. So (S) holds, and Corollary~\ref{cor:coherent-assigned}(c) completes the proof.
\end{proof}
The set $\mathcal K$ may be infinite. Pettigrew's additive Bregman divergences meet the hypotheses whenever they are finite and lower semicontinuous, in particular the Brier score, which he endorses; additivity and the Bregman form are not needed. His Theorem~10.1.1 gives more: improvement under every possible chance function, by a dominator that is not even moderately modest. The proposition gives ordinary chance dominance by an undominated member of $\C$. Either way, the argument requires a normative principle concerning chance dominance, not actual-loss dominance.}

\bibliographystyle{plain}
\bibliography{ldg-references}
\end{document}
